\documentclass[11pt]{article}

\usepackage{amsmath, amsfonts, amsthm, amssymb}


\title{SINDy}

\begin{document}
    \section{Introduction}
        \begin{itemize}
            \item Introduce the problem with discovering physical systems. Can use the same analogy in the
            paper of newton vs kepler or find a different one.
            \item Talk about how using data to find patterns is useful, but we are more interested in finding
            underlying equations that govern the dynmical systems.
            \item We can use discovered equation to make
            predictions in areas where we have never seen data.
        \end{itemize}

    \section{How it works}
    \begin{itemize}
        \item We want to find $f$ for $\dot{x}(t) = f(x(t))$
        \item We start with $n$ observations of our data $x(t)$ at time points $t_0, t_1, t_2,\dots,t_n$
        \item We might also have $m$ variables.
        \item Construct a matrix $X$ of shape $n \times m$, where each row is a snapshot in time, and each
        column is that state variable. We also compute $\dot{X}$
        \item We construct a feature library $\Theta(X)$ of shape $n \times p$ where $p$ is the number of
        candidate functions.
        \item Candidate functions can consist of many things including $x, x^2, \sin(x), \cos(x),\dots$
        \item 
    \end{itemize}
    \section{Finish with examples, ex lotka volterra}
\end{document}